\section*{ЗАКЛЮЧЕНИЕ}
\addcontentsline{toc}{section}{ЗАКЛЮЧЕНИЕ}
\thispagestyle{plain}

В рамках настоящей магистерской работы спроектирована архитектура
и разработан программный сервис EduQuiz для интерактивного
взаимодействия с обучающимися, развёртываемый в инфраструктуре
Российской Федерации и применимый к преподаванию учебных дисциплин
в высшем учебном заведении. Поставленные задачи решены полностью.

Проанализированы распространённые системы интерактивного обучения
(Quizlet, Kahoot!, Mentimeter, Slido, JoyTeka, OnlineTestPad). Показано,
что ни одно из существующих решений не сочетает одновременно работу
в режиме реального времени, расширенный набор игровых механик,
размещение серверной инфраструктуры на территории Российской
Федерации и наличие открытого программного интерфейса.
Рассмотрены теоретические основы геймификации в образовании,
алгоритма интервального повторения SuperMemo SM-2 и психометрики
тестового инструмента средствами классической теории тестов.
На основании проведённого анализа сформулированы семь функциональных
и шесть нефункциональных требований к разрабатываемому сервису.

Обоснован выбор монолитной клиент-серверной архитектуры
с собственным backend-приложением и открытым программным
интерфейсом. Обоснован выбор стека программных средств: Go~1.25 с фреймворком
Gin для серверной части, PostgreSQL~16 и Redis~7 для хранения
данных и кеширования, обратный прокси Caddy~2 с автоматическим
выпуском TLS-сертификата, Flutter~3.35 для кроссплатформенного
клиентского приложения с единой кодовой базой для Web, iOS,
Android и macOS. Архитектура сервиса описана в нотации
контейнерной диаграммы C4.

Разработаны модели данных и алгоритмы сервиса: ролевая модель
пользователей с разграничением доступа на двух уровнях
по принципу наименьших привилегий; модель интерактивной учебной
сессии в виде конечного автомата с поддержкой трёх режимов
проведения (синхронный групповой, индивидуальный с таймером,
индивидуальный соревновательный); модель игровой прогрессии
с формулами расчёта очков опыта и уровней; модель данных в форме
нормализованной реляционной схемы PostgreSQL. Адаптирован алгоритм
индивидуального интервального повторения SuperMemo SM-2 для режима
интерактивной сессии через автоматическое формирование оценки
качества ответа из пары «корректность~--- время», что устраняет
необходимость отдельного интерфейса самооценки и обеспечивает
плавную интеграцию повторения в цикл сессии. Разработана модель
оценки качества банка тестовых вопросов средствами классической
теории тестов: расчёт коэффициента $\alpha$ Кронбаха, индекса
трудности и точечно-бисериальной корреляции пунктов с цветовой
интерпретацией результатов для преподавателя.

Реализованы четыре функциональных модуля клиентского приложения:
модуль аутентификации с поддержкой регистрации, входа и обновления
токенов доступа, разграничением доступа по ролям и комплексом мер
информационной безопасности по перечню OWASP Top 10;
модуль преподавателя с управлением курсами и банками вопросов,
редактором вопросов нескольких типов, импортом банков в форматах
Moodle XML и GIFT и психометрической аналитикой; модуль проведения
интерактивной учебной сессии в трёх режимах с синхронизацией
состояния между клиентами по протоколу WebSocket; модуль
обучающегося с индивидуальным интервальным повторением по алгоритму
SuperMemo SM-2 и профилем игровой прогрессии. Разработаны и приведены
реализации ключевых компонентов серверной части~--- алгоритм SM-2,
формула расчёта очков опыта, психометрический анализатор банка
вопросов. Сервис развёрнут в инфраструктуре Yandex Cloud
с автоматическим деплоем средствами GitHub Actions и доступен
по доменному имени \texttt{eduquizz.ru}.

Методом нагрузочного тестирования собственным программным
инструментом установлено, что разработанная архитектура в одиночной
серверной конфигурации устойчиво обслуживает до 500 одновременных
обучающихся в одной интерактивной сессии при перцентиле $p_{99}$
латентности критического пути обработки запросов не выше 233~мс
и нулевой доле ошибок класса 5xx. Нефункциональное требование
к производительности выполнено.

Разработана методика квазиэкспериментального педагогического
исследования по схеме pre-/post-/retention-тестирования
и реализован программный пайплайн обработки результатов
на языке Python с применением критериев Шапиро\textendash Уилка
и Манна\textendash Уитни, оценки размера эффекта по коэффициенту
Коэна, поправки Бонферрони, расчёта коэффициента $\alpha$ Кронбаха
и точечно-бисериальной корреляции пунктов. Корректность пайплайна
проверена на модельных данных; апробация педагогической методики
на реальной группе обучающихся курса \textit{«Надёжность»} кафедры
автоматизированных систем управления факультета автоматики
и вычислительной техники РГУ Нефти и Газа (НИУ) им. И.М. Губкина
запланирована в текущем семестре.

Сформулированы перспективы дальнейшей научно-исследовательской
работы~--- задел для последующих этапов исследования: полная
двусторонняя интеграция с системой дистанционного обучения
университета; промышленный мониторинг и резервирование критичных
узлов; разработка модели латентного уровня подготовки обучающегося
на основе теории отклика на пункт и реализация адаптивного режима
тестирования; разработка модели оценки валидности атрибуции
учебного ответа в условиях распространения общедоступных больших
языковых моделей и интегральной оценки знаний, учитывающей
одновременно измерительную точность тестового инструмента
и валидность атрибуции попытки.

Таким образом, в работе реализован программный сервис интерактивного
взаимодействия с обучающимися EduQuiz, обеспечивающий проведение
учебных сессий в режиме реального времени, индивидуальное повторение
учебного материала и функционирование в инфраструктуре российского
высшего учебного заведения. Полученное решение создаёт основу
для применения сервиса в учебном процессе кафедры автоматизированных
систем управления РГУ Нефти и Газа (НИУ) имени И.М. Губкина
и для последующего развития комплекса в рамках самостоятельных
научно-исследовательских направлений.

\newpage
